\documentclass[11pt]{article}

% Main report for the GPU roofline profiler. Every number this document shows is
% generated by the analysis pipeline and pulled in from report/figures and
% report/tables. Where a generated artifact does not exist yet, a guard prints a
% visible "pending" note so the document always compiles and never shows a stale
% or invented figure. No performance number is ever typed by hand.

\usepackage[margin=1in]{geometry}
\usepackage{amsmath,amssymb}
\usepackage{booktabs}
\usepackage{longtable}
\usepackage{graphicx}
\usepackage{siunitx}
\usepackage{xcolor}
\usepackage{caption}
\usepackage{float}
\usepackage[hidelinks]{hyperref}

\graphicspath{{figures/}}

% Pull in a generated fragment (table or similar) if it exists, otherwise print
% a visible placeholder. This is what lets the report compile before and after
% the pipeline runs, with no hand editing in between.
\newcommand{\inputgenerated}[1]{%
  \IfFileExists{#1}{\input{#1}}{%
    \textcolor{red}{[pending pipeline output: \texttt{\detokenize{#1}}]}}}

% Include a generated figure if present, otherwise a visible placeholder box.
\newcommand{\includegenerated}[2][width=\linewidth]{%
  \IfFileExists{figures/#2}{\includegraphics[#1]{#2}}{%
    \fbox{\parbox{0.9\linewidth}{\centering\vspace{2em}%
      \textcolor{red}{figure pending: \texttt{figures/\detokenize{#2}}}\vspace{2em}}}}}

% A generated table that is wider than the text block. Scaled down to fit rather
% than allowed to run off the right of the page, which is what the Nsight
% counter table with its seven columns would otherwise do.
\newcommand{\inputgeneratedwide}[1]{%
  \IfFileExists{#1}{\begin{center}\resizebox{\textwidth}{!}{\input{#1}}\end{center}}{%
    \textcolor{red}{[pending pipeline output: \texttt{\detokenize{#1}}]}}}

% A generated table set smaller. The appendix tables carry raw Nsight metric
% names, some of them 49 characters with no break point, which overflow the text
% block at the body font size. These cannot be scaled with resizebox because
% they are longtables and span pages, so the font is reduced instead.
\newcommand{\inputgeneratedsmall}[1]{%
  \IfFileExists{#1}{{\footnotesize\input{#1}}}{%
    \textcolor{red}{[pending pipeline output: \texttt{\detokenize{#1}}]}}}

% Title metadata, which hyperref also puts in the PDF properties. The typeset
% title page below is built by hand rather than by \maketitle, because the
% default stacks title, author, and date together at the top of the page.
\title{GPU Roofline Profiler}
\author{Olajide Badejo}
\date{\today}

\begin{document}

% Title vertically and horizontally centred, with the byline pushed to the foot
% of the page. The two \fill glues take up all the slack: equal amounts above
% and below the title centre it, and whatever is left drops the author and date
% to the bottom. \end{titlepage} starts the next page itself, so no \clearpage
% follows it; adding one would insert a blank page.
\begin{titlepage}
  \centering
  \vspace*{\fill}
  {\LARGE\bfseries GPU Roofline Profiler\par}
  \vspace*{\fill}
  {\large Olajide Badejo\par}
  \vspace{0.9em}
  {\large \today\par}
  \vspace{1.5cm}
\end{titlepage}

\begin{abstract}
I measure the compute and memory ceilings of a single RTX 5070, place a ladder
of hand written CUDA kernels from SAXPY to a register blocked GEMM on a roofline
against those ceilings, and explain where each kernel lands using Nsight Compute
counters. Both the theoretical ceilings from device properties and the ceilings
I actually reach with an FMA microbenchmark and a large device copy are plotted,
and the gap between them is discussed rather than hidden. Every figure and table
here is regenerated from raw results by one command; nothing is a spec sheet
number presented as measured.
\end{abstract}
\clearpage

\tableofcontents
\clearpage

\section{Introduction}

This report is about one graphics card and what it can actually do. The card is
an RTX 5070, a Blackwell part with 12\,GB of GDDR7 that also happens to be
driving the display on the machine I built this on. The vendor quotes a peak
around \SI{31}{TFLOPS} of FP32 and a memory bandwidth around \SI{672}{GB/s}. Those
two numbers are the headline of any spec sheet, and on their own they say almost
nothing about how fast a given kernel will run, because most kernels are limited
not by raw compute but by how much data they have to move to feed that compute.

The roofline model exists to make that trade off visible. It plots attainable
performance against arithmetic intensity, the ratio of floating point work to
bytes moved, and bounds every kernel under two ceilings at once: a flat compute
ceiling and a sloped memory ceiling. Where a kernel sits under that roof tells
you which resource it is starved on, and therefore what would make it faster.

The contribution here is not another fast GEMM. It is a ladder of kernels that
walks the arithmetic intensity axis on purpose, from a SAXPY that touches memory
and does almost no arithmetic, through a transpose that does no floating point
work at all, up to a register blocked GEMM that reuses each loaded value many
times. For every rung I measure where it lands, and I explain that position with
the specific hardware counters behind it, not with a generic statement about
what tiling usually does. The two ceilings I draw are both real: one derived
from the device properties, one measured by pushing the card as hard as I can
with a microbenchmark, and I am honest about the distance between them.

\section{The roofline model}

The roofline model of Williams, Waterman, and Patterson~\cite{williams2009roofline}
bounds the performance a kernel can attain on a machine by the smaller of two
limits. Let $P_{\max}$ be the peak floating point rate in FLOP/s and $B$ the peak
memory bandwidth in byte/s. For a kernel with arithmetic intensity $I$, measured
in FLOP per byte, the attainable performance is

\begin{equation}
  P(I) = \min\!\left(P_{\max},\; I \cdot B\right).
\end{equation}

The first term is a horizontal line, the compute ceiling: no amount of memory
bandwidth lets a kernel exceed the rate at which the arithmetic units retire
operations. The second term is a line through the origin with slope $B$, the
memory ceiling: a kernel of intensity $I$ cannot run faster than the rate at
which memory can deliver its operands. Plotted on log-log axes, the two combine
into the characteristic roof shape.

The two ceilings meet at the ridge point,

\begin{equation}
  I_{\text{ridge}} = \frac{P_{\max}}{B}.
\end{equation}

A kernel whose intensity is below $I_{\text{ridge}}$ is memory bound: it sits
under the sloped part of the roof, and the way to make it faster is to move less
data or move it more efficiently. A kernel above the ridge is compute bound: it
sits under the flat part, and only more arithmetic throughput helps. The ridge
point is a property of the machine, not the kernel, so the same intensity can be
memory bound on one card and compute bound on another.

Arithmetic intensity is where the honesty of the analysis lives. The theoretical
intensity of a kernel counts the minimum bytes a perfect implementation would
move: read each input once, write each output once. The measured intensity uses
the bytes the kernel actually moved through DRAM, which Nsight Compute reports,
and which include cache misses, register spills, and redundant loads that a real
implementation suffers. The gap between the two is most of the interesting story
in this report, so every point on my rooflines is labeled with which byte count
produced its intensity.

\section{Methodology}

\subsection{Timing}

All kernel timing uses CUDA events, not host wall clock, so the measurement sees
device time and is not polluted by launch overhead or host scheduling. To keep
the per launch overhead from dominating on the cheap kernels, I place many
launches between a single pair of events and divide by the launch count rather
than synchronizing once per launch. Before any timed run I discard a batch of
warmup iterations, which lets the boost clock ramp and the caches warm so the
timed region measures steady state rather than the cold transient.

I then time at least ten batches and report the mean, median, and standard
deviation. This card boosts its clock opportunistically and shares silicon with
the display, so run to run variance is real and worth showing rather than hiding
behind a single number. The standard deviation column in every timing table is
there for exactly that reason.

\subsection{Counting floating point operations}

Each kernel's FLOP count comes from its actual arithmetic, not an estimate. A
GEMM of dimensions $M \times N \times K$ does $2MNK$ floating point operations,
counting each multiply and add separately. A GEMV of an $M \times N$ matrix does
$2MN$. A SAXPY over $N$ elements does $2N$, one multiply and one add per element.
A transpose does zero floating point work, which is the whole point of including
it: it isolates the memory system with no arithmetic to hide behind.

\subsection{Counting bytes}

I report two byte counts wherever both are available. The theoretical count is
the minimum traffic a perfect implementation would generate: read each input
once, write each output once, in the kernel's precision. The measured count is
the DRAM read and write bytes Nsight Compute observed. Arithmetic intensity is
computed from the measured bytes where the profiler ran, and falls back to the
theoretical bytes elsewhere, always labeled so the two are never confused.

\subsection{Theoretical versus empirical peak}

The theoretical ceilings come from the device properties queried at runtime:
the FP32 peak from the SM count, the FP32 lanes per SM, and the boost clock,
counting a fused multiply add as two operations; the bandwidth peak from the bus
width and the effective memory transfer rate. The empirical ceilings come from
pushing the hardware directly: a register resident FMA loop for compute, and a
large device to device copy for bandwidth. A card almost never reaches its
theoretical peak, and reporting the measured ceiling next to the theoretical one
is what keeps the roofline honest about what this specific machine delivers.

\section{Kernel implementations}

The suite is a ladder along the arithmetic intensity axis. Each rung is a kernel
whose position on the roofline is meant to teach something specific.

\subsection{SAXPY}

$y \leftarrow a x + y$ over long vectors. Two floating point operations per
element against three memory accesses (two reads, one write), so its intensity is
tiny and it lives at the far memory bound left edge of the roofline. It is the
anchor: if SAXPY does not reach close to the measured bandwidth ceiling, the
harness or the copy microbenchmark is wrong before any GEMM is even considered.

\subsection{Reduction}

A shared memory tree reduction that sums a large array. It exercises the
synchronization and shared memory path rather than raw streaming, and it is
tested across non power of two sizes because the tail handling is where reduction
kernels usually break.

\subsection{Transpose, naive and tiled}

The transpose does no floating point work, which isolates coalescing and shared
memory bank conflicts with nothing else in the way. The naive version reads the
input coalesced but writes the transpose with a large stride, so the writes are
uncoalesced. The tiled version stages a tile in shared memory so both the global
read and the global write are coalesced. The tile is declared one column wider
than it is tall, \texttt{tile[TILE\_DIM][TILE\_DIM+1]}, so that the transposed
access pattern to shared memory does not map every thread in a warp onto the same
bank. That single pad word removes the bank conflicts on the transposed access,
and the naive versus tiled delta is the cleanest demonstration in the whole
suite.

\subsection{GEMV}

Matrix times vector. Each matrix element is read once and touched by one
multiply add, so reuse is low and the kernel sits in the middle of the intensity
axis, between the streaming kernels and the GEMM ladder.

\subsection{The GEMM ladder}

General matrix multiply is where intensity and achieved performance climb
together as reuse improves, and it is the core of the report.

\begin{description}
  \item[Naive.] Every thread reads its row and column straight from global
    memory. Correct, and slow, because each value is reloaded from DRAM many
    times. This is the memory bound bottom of the GEMM ladder.
  \item[Shared memory tiled.] Threads cooperate to stage tiles of $A$ and $B$
    into shared memory once, then reuse them across the tile. Reuse rises and the
    kernel climbs the roof.
  \item[Register blocked.] Each thread computes a small output tile rather than a
    single element, holding partial sums in registers. This raises the reuse per
    shared memory load, trading register pressure for far fewer shared accesses,
    and it is usually the step with the largest jump in achieved performance.
  \item[Vectorized.] The register blocked kernel with \texttt{float4} loads and
    stores, so each memory instruction moves four values and the load and store
    issue pressure drops.
  \item[cuBLAS SGEMM.] The vendor library, plotted as an honest ceiling for a
    custom kernel to be measured against. It is labeled as a library, not as one
    of my kernels, because it is not a fair comparison of hand written code and
    it would be dishonest to present it as one.
\end{description}

A WMMA tensor core GEMM is a stretch goal on its own roofline panel, since its
compute ceiling is a different number from the FP32 CUDA core ceiling and the two
should not share an axis.

\section{Experimental setup}

The measurements in this report come from one machine, described in the
environment appendix, which is generated from the run manifests rather than typed
here. The GPU is an RTX 5070 on Windows 11, and it is also the display adapter,
so the benchmark driver checks free memory before each allocation and skips any
configuration that would not leave comfortable headroom for the desktop.

Every pipeline run writes a timestamped directory under \texttt{results/raw} with
a manifest recording the resolved sweep configuration, the driver and CUDA
versions, an \texttt{nvidia-smi} snapshot, and the exact Nsight tool versions.
Results without a manifest are treated as not existing. The sweep parameters,
the sizes, the tile dimensions, and the repeat counts, all live in
\texttt{configs/sweep.yaml}; nothing about the experiment is hard coded in the
source.

The kernels profiled in detail with Nsight Compute are a curated subset of the
timing sweep, each at two or three representative sizes, because the profiler
replays each kernel many times per metric set and profiling the full sweep would
take hours to no additional benefit.

\inputgenerated{tables/environment_summary.tex}

\section{Results}

\subsection{Ceilings}

The theoretical and measured ceilings for this card are collected in the table
below. The theoretical numbers are derived from device properties; the measured
numbers come from the FMA microbenchmark and the large device to device copy.

\inputgenerated{tables/ceilings.tex}

\subsection{The roofline}

Figure~\ref{fig:roofline-main} places every kernel configuration on the roofline
against both ceilings. Points whose intensity was computed from measured DRAM
bytes are drawn differently from points that fall back to a theoretical byte
count, so the reader always knows which is which.

% [H] pins the figure exactly here. Left as a floating [htbp] it was pushed to a
% page of its own and left the remainder of the page blank behind it.
\begin{figure}[H]
  \centering
  \includegenerated{roofline_main.pdf}
  \caption{Kernel ladder on the RTX 5070 roofline, theoretical and measured
    ceilings shown together. One label per kernel family gives the largest
    problem size that family was run at.}
  \label{fig:roofline-main}
\end{figure}

\subsection{Nsight Compute counters}

The counters that explain each kernel's position, achieved occupancy, DRAM
throughput, shared memory bank conflicts, and coalescing measured as sectors per
request, are collected for the curated subset below. A perfectly coalesced warp
moves 4 sectors per request for a 4 byte access and 16 for a \texttt{float4}, so
the ideal depends on the access width and is not a single number.

\inputgeneratedwide{tables/ncu_counters.tex}

\subsection{Timing summary}

Every configuration in the sweep. The table runs over several pages.

\inputgenerated{tables/timing_summary.tex}

\section{Discussion and bottleneck analysis}

This is where the project earns its title. A roofline plot on its own only shows
that a kernel is slow; it does not say why. The Nsight Compute counters do, and
each kernel's position below is tied to the specific counter that explains it
rather than to a generic claim about what an optimization usually achieves. The
prose in this section is written against the collected counters once the profiling
pass has run on this machine; until then the structure below records exactly
which counter is expected to explain each rung, so the analysis is a fill in of
measured evidence and not a rewrite.

\subsection{The streaming kernels against the bandwidth ceiling}

SAXPY and the reduction should sit close under the measured bandwidth ceiling.
The counter that settles it is the DRAM throughput as a fraction of peak: if
SAXPY reaches most of the measured copy bandwidth, the harness is sound and the
memory bound edge is anchored. Any large shortfall points at launch overhead or
occupancy rather than bandwidth, and the occupancy counter distinguishes them.

\subsection{Transpose: coalescing and bank conflicts}

The naive and tiled transpose is the cleanest counter story in the suite. The
naive kernel's uncoalesced writes show up as a low global store efficiency; the
tiled kernel fixes that and the efficiency counter jumps. The padded tile then
shows up in the shared memory bank conflict counter falling to near zero, where
an unpadded tile would show heavy conflicts on the transposed access. The two
counters, store efficiency and bank conflicts, are exactly the two levers the
tiled transpose pulls.

\subsection{The GEMM ladder: reuse and occupancy}

Up the GEMM ladder, arithmetic intensity and achieved performance should climb
together, and the counters say why at each step. The naive kernel is memory
bound with low intensity. Shared memory tiling raises the intensity by cutting
DRAM traffic, visible as the measured intensity moving right toward the
theoretical value. Register blocking raises reuse further but spends registers,
so the occupancy counter is where its cost shows up: the interesting question is
whether the reuse gain outweighs the occupancy it costs, and the counters answer
it directly. Vectorized loads reduce the memory instruction count, and cuBLAS
sits above the custom kernels as the vendor reference.

\subsection{Thermals and clocks}

Because this is a boost clocked consumer card also driving a display, the NVML
log recorded alongside every run is what shows whether a result was shaped by
temperature or clock drift rather than by kernel design. Where a kernel's timing
variance is large, the joined power and clock samples say whether the card was
throttling or simply boosting unevenly, and that distinction is called out here
wherever the NVML trace makes it.

\section{Conclusion and future work}

The suite walks a kernel ladder from a memory bound SAXPY to a compute bound
register blocked GEMM, places each rung on a roofline against both the
theoretical and the measured ceilings for this RTX 5070, and ties each position
back to the Nsight Compute counter that explains it. The value is not any single
fast kernel but the counter backed account of why each kernel lands where it
does, and the honest gap between the vendor spec sheet and what the card actually
delivers.

Several extensions would sharpen the picture. A WMMA tensor core GEMM on its own
roofline panel would add a second, higher compute ceiling and show how far the
FP32 CUDA core story is from the tensor core story on the same silicon. A cache
aware roofline with a second sloped ceiling for L2 bandwidth would separate
kernels that are DRAM bound from kernels that are L2 bound, which the single
DRAM ceiling here cannot distinguish. Sweeping the power and clock limits and
replotting would turn the NVML trace from a diagnostic into a first class axis,
showing how the ceilings themselves move as the card is throttled. Each of these
is additive: the harness, the analysis package, and the report pipeline are built
to absorb a new kernel or a new ceiling and regenerate every figure with one
command.


\appendix
\section{Environment and raw tables}

The machine environment below is generated at report build time from the run
manifests, never typed from memory, so it always describes the machine that
produced the results in this document.

\inputgenerated{tables/environment_full.tex}

\subsection{Full timing table}

\inputgenerated{tables/timing_full.tex}

\subsection{Full Nsight Compute metric table}

\inputgenerated{tables/ncu_full.tex}


\bibliographystyle{plain}
\bibliography{references}

\end{document}
